\section{Deviations from the frozen protocol}
\label{app:deviations}

The protocol (hypotheses, the 35-test family, configurations, data manifests and code) was frozen on 2026-10-08 at
07:45 UTC and recorded as SHA-256 hashes. While the local decoder runs were in progress, three code audits
were carried out; their findings led to one amendment, frozen at 10:58 UTC the same day. By then, test predictions of
Qwen3.5-4B and Tev1-4B had been produced, and the run log had printed their test accuracies. No test of the
35-test family, no shift analysis and no DeBERTa or Jev run had been carried out. The amendment did not change the code
used by the local decoder models, so all decoder runs, before and after it, use identical code; the DeBERTa and Jev
runs and every analysis use the amended version.
\paragraph{Chronology (UTC, from the run records).}
\begin{itemize}
  \item 2026-10-06: label texts frozen. 2026-10-07: data re-processed after a cleaning fix; candidate sets rebuilt as
  fixed groups after the leakage audit; pipeline checks, some on test splits (see the disclosure below).
  \item 2026-10-08 07:45: protocol v1.1 frozen (hash file).
  \item 2026-10-08 09:14--10:45: all Qwen3.5-4B and Tev1-4B predictions (development and test); 10:45: first
  Qwen3.5-9B run started.
  \item 2026-10-08 10:58: amendment v1.1.1 frozen; no test of the 35-test family and no shift analysis had been run.
  \item 2026-10-08 10:45--16:56: remaining local decoder predictions; latency timing completed on 2026-10-09.
  2026-10-09 to 10-10: DeBERTa-v3-large. 2026-10-13 05:23: TF-IDF.
  \item 2026-10-13 05:42: no-shift tables and P1/P2 for all systems except Jev; 05:47--07:50: all Jev runs;
  07:50--08:19: Jev's no-shift table and the descriptive shift experiment; 2026-10-13/14: PA1--PA3; 2026-10-14: Holm
  over all 35 tests.
  \item 2026-10-14, after all confirmatory results: prior-set sensitivity of PA1--PA3, HWU64 without the MASSIVE
  overlap, descriptive change in ECE and recalibration arms for all systems.
  \item 2026-10-15: distance of the implicit priors from the true class distribution, ECE range of the repeated Jev
  runs, excess-risk, Brier and NLL tables, TF-IDF grid extension to $10^5$ and $10^6$, average of PA1--PA3 over 2{,}000 priors, EM without
  shift, FPB confusion matrices, rank correlations between metrics. All are listed below and none changes a
  confirmatory result.
\end{itemize}

\paragraph{Amendment v1.1.1.}
\begin{itemize}
  \item \textbf{No-shift control.} The shift effect on ECE (the descriptive shift experiment and PA1) is measured against a no-shift control
  drawn with the same with-replacement resampler at the empirical class prior. The original protocol compared against
  the unresampled test set, which would have attributed resampling bias in top-label ECE to the shift.
  \item \textbf{Bootstrap size for PA1--PA3.} $B$ was raised from 1{,}000 to 10{,}000. With $B=1{,}000$ the smallest
  non-zero two-sided $p$ ($2/B=0.002$) is above Holm's first threshold ($0.05/35\approx0.0014$).
  \item \textbf{Failed-run rule for DeBERTa.} The majority-class threshold is computed within each example's candidate
  set, on the same scale as the model's accuracy (previously the two were on different scales for CLINC150 and HWU64).
  No full-data run failed under either rule.
  \item \textbf{Jev transport errors.} Network errors are never classified as refusals. No Jev request failed or was
  refused in the final runs.
  \item \textbf{No-shift metrics.} The specified treatment of refused requests (counted as wrong and abstained) was
  implemented; no request was refused. The confirmatory scope of H3 was narrowed in wording to plain EM on Jev, and
  the AURC definition was stated exactly.
\end{itemize}
\textbf{Analyses added after the results were known.} The following analyses were added after the confirmatory results
had been computed, in response to reviews of the draft; all are labelled exploratory or descriptive and none
changes a confirmatory result: the prior-set sensitivity of PA1--PA3 (\hyperref[tab:pasens]{Table~\ref*{tab:pasens}}) and their average over
2{,}000 target priors (\hyperref[tab:papop]{Table~\ref*{tab:papop}}); the HWU64 analysis without the MASSIVE overlap (required by the protocol
but run only at this stage; \hyperref[tab:hwu64]{Table~\ref*{tab:hwu64}}); the descriptive shift-experiment summaries for all systems
(\hyperref[tab:exp2delta]{Tables~\ref*{tab:exp2delta}}--\nobreak\ref{tab:nll}); the distance of each implicit prior from the true class distribution
(\hyperref[tab:implicitprior]{Table~\ref*{tab:implicitprior}}); the effect of EM without shift (\hyperref[tab:emnoshift]{Table~\ref*{tab:emnoshift}}); the FPB confusion matrices
of the Nimble pair (\hyperref[tab:fpbconf]{Table~\ref*{tab:fpbconf}}); the ECE range of the repeated Jev runs
(\hyperref[tab:jevrepeat]{Table~\ref*{tab:jevrepeat}}); the rank correlations between metrics
(\hyperref[tab:exp1full]{Table~\ref*{tab:exp1full}}); a check that the PA estimates are not an artefact of
resampling twice (PA1--PA3 recomputed without the bootstrap base on FPB, HateXplain and TREC-6; the estimates agree);
and the extension of the TF-IDF regularisation grid to $10^5$ and $10^6$
(\hyperref[sec:protocol]{Section~\ref*{sec:protocol}}).

\textbf{Disclosure.} During pipeline development, before the freeze, checks of the code were run on \emph{test}
splits: Jev on the full FPB test split and on 5 to 10 test examples of CLINC150, FPB and HateXplain; Qwen3.5-4B and
Tev1-4B on 5 FPB test examples each; a DeBERTa-v3-large model trained on 100 examples, on the FPB test split; TF-IDF
on all five test splits in three pipeline runs, whose test scores were printed; and early versions of PA1--PA3 and
the descriptive shift experiment on FPB. Two design changes followed these checks: fixed target priors for PA1--PA3,
because of the width of the intervals under the earlier design, and the extension of the TF-IDF grid from $10^3$ to
$10^4$, because the value selected on the development set lay at the edge of the grid. Neither change was based on
the direction or size of a test estimate. From then on, pipeline checks used only the development split, and all reported test results come from the
frozen protocol.

\FloatBarrier
\section{Additional results}
\label{app:results}

\begin{table}[htbp]
  \centering
  \caption{\textbf{Fourteen of the 35 confirmatory tests are significant after Holm correction.} Paired differences in percentage points (TDM $-$ base for P1/P2; Jev $-$
  DeBERTa-v3-large for PA1/PA2; EM $-$ none for PA3) with unadjusted 95\% bootstrap intervals, raw two-sided $p$ and
  Holm-adjusted $p$ over the 35 tests.}
  \label{tab:alltests}
  \scriptsize
  \begin{tabular}{llrcrr}
\toprule
Test & Dataset & $\Delta$ (points) & 95\% CI & $p$ & Holm $p$ \\
\midrule
\rowcolor{rowhl} P1-F1 & CLINC150 & +2.48 & [+1.94, +3.10] & $<$0.0002 & $<$0.007 \\
\rowcolor{rowhl} P1-F1 & HWU64 & +2.84 & [+1.10, +4.78] & 0.002 & 0.040 \\
\rowcolor{rowhl} P1-F1 & FPB & +9.99 & [+7.35, +12.69] & $<$0.0002 & $<$0.007 \\
P1-F1 & HateXplain & +1.69 & [$-$0.16, +3.52] & 0.077 & 1.000 \\
P1-F1 & TREC-6 & +0.65 & [$-$2.92, +4.44] & 0.724 & 1.000 \\
\midrule
P1-ECE & CLINC150 & $-$0.34 & [$-$0.84, +0.21] & 0.255 & 1.000 \\
P1-ECE & HWU64 & $-$1.87 & [$-$3.29, $-$0.19] & 0.028 & 0.479 \\
\rowcolor{rowhl} P1-ECE & FPB & $-$13.19 & [$-$15.75, $-$10.39] & $<$0.0002 & $<$0.007 \\
\rowcolor{rowhl} P1-ECE & HateXplain & $-$3.23 & [$-$5.19, $-$1.37] & 0.001 & 0.029 \\
P1-ECE & TREC-6 & +0.12 & [$-$2.97, +2.37] & 0.796 & 1.000 \\
\midrule
\rowcolor{rowhl} P2-F1 & CLINC150 & +0.69 & [+0.23, +1.17] & 0.001 & 0.029 \\
P2-F1 & HWU64 & +1.14 & [$-$0.64, +3.03] & 0.195 & 1.000 \\
\rowcolor{rowhl} P2-F1 & FPB & $-$11.97 & [$-$14.61, $-$9.28] & $<$0.0002 & $<$0.007 \\
P2-F1 & HateXplain & $-$1.66 & [$-$3.23, $-$0.07] & 0.039 & 0.630 \\
P2-F1 & TREC-6 & $-$3.11 & [$-$6.99, +0.69] & 0.100 & 1.000 \\
\midrule
\rowcolor{rowhl} P2-ECE & CLINC150 & $-$1.22 & [$-$1.52, $-$0.63] & $<$0.0002 & $<$0.007 \\
\rowcolor{rowhl} P2-ECE & HWU64 & $-$4.22 & [$-$5.59, $-$2.50] & $<$0.0002 & $<$0.007 \\
\rowcolor{rowhl} P2-ECE & FPB & +9.56 & [+6.77, +12.05] & $<$0.0002 & $<$0.007 \\
P2-ECE & HateXplain & $-$0.36 & [$-$1.94, +1.21] & 0.636 & 1.000 \\
P2-ECE & TREC-6 & $-$1.84 & [$-$3.67, +1.63] & 0.440 & 1.000 \\
\midrule
PA1 & CLINC150 & +0.09 & [$-$0.02, +0.20] & 0.104 & 1.000 \\
PA1 & HWU64 & +0.11 & [$-$0.30, +0.50] & 0.592 & 1.000 \\
\rowcolor{rowhl} PA1 & FPB & +1.71 & [+0.98, +2.49] & $<$0.0002 & $<$0.007 \\
PA1 & HateXplain & $-$0.78 & [$-$1.29, $-$0.26] & 0.004 & 0.080 \\
PA1 & TREC-6 & +1.01 & [+0.26, +1.85] & 0.007 & 0.136 \\
\midrule
PA2 & CLINC150 & 0.00 & [$-$0.04, +0.04] & 0.938 & 1.000 \\
PA2 & HWU64 & +0.01 & [$-$0.21, +0.23] & 0.969 & 1.000 \\
\rowcolor{rowhl} PA2 & FPB & +2.49 & [+1.83, +3.16] & $<$0.0002 & $<$0.007 \\
PA2 & HateXplain & $-$0.55 & [$-$0.98, $-$0.14] & 0.010 & 0.187 \\
PA2 & TREC-6 & +0.55 & [$-$0.06, +1.30] & 0.086 & 1.000 \\
\midrule
PA3 & CLINC150 & $-$0.12 & [$-$0.29, +0.03] & 0.122 & 1.000 \\
PA3 & HWU64 & +0.12 & [$-$0.33, +0.54] & 0.546 & 1.000 \\
\rowcolor{rowhl} PA3 & FPB & $-$2.65 & [$-$3.74, $-$1.61] & $<$0.0002 & $<$0.007 \\
\rowcolor{rowhl} PA3 & HateXplain & $-$1.75 & [$-$2.71, $-$0.82] & $<$0.0002 & $<$0.007 \\
PA3 & TREC-6 & $-$1.06 & [$-$1.96, $-$0.28] & 0.007 & 0.136 \\
\bottomrule
\end{tabular}

\end{table}

\begin{table}[htbp]
  \centering
  \caption{\textbf{Accuracy and MCC rank the seven analysed systems almost as Macro-F1 does (Spearman $\rho\ge0.75$ on
  every dataset; descriptive), whereas NLL, which penalises confident errors, ranks them differently on HWU64 and
  HateXplain ($\rho=0.29$ and $0.32$).} Test-set metrics for all systems, no shift. $^\ast$Exploratory: a later Nimble
  release whose training data include public TREC questions, which may explain its TREC-6 score; it is not part of any
  confirmatory test or of the rank correlations. Acc@80: accuracy on the 80\% most confident examples.}
  \label{tab:exp1full}
  \scriptsize
  \setlength{\tabcolsep}{3pt}
  \begin{tabular}{llrrrrrrrr}
\toprule
Dataset & System & Acc$\uparrow$ & F1$\uparrow$ & MCC$\uparrow$ & ECE$\downarrow$ & Brier$\downarrow$ & NLL$\downarrow$ & AURC$\downarrow$ & Acc@80$\uparrow$ \\
\midrule
CLINC150 & Jev & 0.978 & 0.978 & 0.978 & 0.007 & 0.032 & 0.152 & 0.004 & 0.997 \\
 & Tev1-4B & 0.964 & 0.964 & 0.963 & 0.008 & 0.055 & 0.130 & 0.003 & 0.997 \\
 & Nimble-9B & 0.962 & 0.962 & 0.961 & 0.008 & 0.060 & 0.140 & 0.004 & 0.996 \\
 & Nimble-9B (main)$^\ast$ & 0.969 & 0.969 & 0.968 & 0.010 & 0.050 & 0.123 & 0.002 & 0.997 \\
 & Qwen3.5-4B & 0.939 & 0.939 & 0.939 & 0.011 & 0.089 & 0.233 & 0.007 & 0.994 \\
 & Qwen3.5-9B & 0.955 & 0.955 & 0.955 & 0.020 & 0.067 & 0.176 & 0.005 & 0.995 \\
 & DeBERTa-v3-large & 0.980 & 0.980 & 0.980 & 0.013 & 0.031 & 0.105 & 0.002 & 0.997 \\
 & TF-IDF + LR & 0.950 & 0.950 & 0.949 & 0.015 & 0.077 & 0.189 & 0.005 & 0.993 \\
\midrule
HWU64 & Jev & 0.911 & 0.908 & 0.910 & 0.036 & 0.138 & 0.640 & 0.024 & 0.976 \\
 & Tev1-4B & 0.891 & 0.895 & 0.890 & 0.040 & 0.168 & 0.423 & 0.021 & 0.962 \\
 & Nimble-9B & 0.897 & 0.896 & 0.896 & 0.035 & 0.165 & 0.436 & 0.026 & 0.956 \\
 & Nimble-9B (main)$^\ast$ & 0.896 & 0.896 & 0.895 & 0.042 & 0.156 & 0.360 & 0.017 & 0.974 \\
 & Qwen3.5-4B & 0.862 & 0.866 & 0.860 & 0.058 & 0.216 & 0.596 & 0.029 & 0.944 \\
 & Qwen3.5-9B & 0.880 & 0.884 & 0.879 & 0.078 & 0.200 & 0.616 & 0.028 & 0.956 \\
 & DeBERTa-v3-large & 0.943 & 0.939 & 0.942 & 0.030 & 0.094 & 0.269 & 0.009 & 0.984 \\
 & TF-IDF + LR & 0.901 & 0.900 & 0.900 & 0.098 & 0.177 & 0.444 & 0.019 & 0.966 \\
\midrule
FPB & Jev & 0.744 & 0.752 & 0.608 & 0.095 & 0.373 & 0.745 & 0.120 & 0.795 \\
 & Tev1-4B & 0.715 & 0.735 & 0.597 & 0.072 & 0.378 & 0.582 & 0.141 & 0.751 \\
 & Nimble-9B & 0.608 & 0.629 & 0.508 & 0.263 & 0.612 & 1.138 & 0.205 & 0.644 \\
 & Nimble-9B (main)$^\ast$ & 0.652 & 0.683 & 0.542 & 0.167 & 0.509 & 0.849 & 0.194 & 0.688 \\
 & Qwen3.5-4B & 0.597 & 0.635 & 0.499 & 0.204 & 0.585 & 1.107 & 0.252 & 0.640 \\
 & Qwen3.5-9B & 0.743 & 0.749 & 0.607 & 0.168 & 0.403 & 0.821 & 0.115 & 0.818 \\
 & DeBERTa-v3-large & 0.875 & 0.870 & 0.776 & 0.111 & 0.234 & 0.634 & 0.045 & 0.920 \\
 & TF-IDF + LR & 0.753 & 0.693 & 0.538 & 0.167 & 0.401 & 0.990 & 0.123 & 0.817 \\
\midrule
HateXplain & Jev & 0.535 & 0.514 & 0.348 & 0.279 & 0.704 & 2.992 & 0.330 & 0.574 \\
 & Tev1-4B & 0.464 & 0.458 & 0.235 & 0.214 & 0.702 & 1.297 & 0.435 & 0.501 \\
 & Nimble-9B & 0.490 & 0.470 & 0.284 & 0.336 & 0.800 & 1.735 & 0.416 & 0.521 \\
 & Nimble-9B (main)$^\ast$ & 0.504 & 0.495 & 0.286 & 0.266 & 0.721 & 1.428 & 0.397 & 0.532 \\
 & Qwen3.5-4B & 0.460 & 0.441 & 0.237 & 0.247 & 0.740 & 1.415 & 0.457 & 0.493 \\
 & Qwen3.5-9B & 0.501 & 0.487 & 0.289 & 0.340 & 0.803 & 1.764 & 0.409 & 0.522 \\
 & DeBERTa-v3-large & 0.710 & 0.703 & 0.563 & 0.141 & 0.433 & 0.769 & 0.161 & 0.774 \\
 & TF-IDF + LR & 0.663 & 0.645 & 0.486 & 0.069 & 0.461 & 0.814 & 0.200 & 0.720 \\
\midrule
TREC-6 & Jev & 0.930 & 0.922 & 0.912 & 0.042 & 0.116 & 0.342 & 0.016 & 0.970 \\
 & Tev1-4B & 0.884 & 0.889 & 0.858 & 0.038 & 0.190 & 0.361 & 0.042 & 0.927 \\
 & Nimble-9B & 0.848 & 0.872 & 0.814 & 0.043 & 0.207 & 0.370 & 0.037 & 0.925 \\
 & Nimble-9B (main)$^\ast$ & 0.968 & 0.963 & 0.960 & 0.017 & 0.056 & 0.129 & 0.006 & 0.990 \\
 & Qwen3.5-4B & 0.880 & 0.882 & 0.853 & 0.037 & 0.177 & 0.363 & 0.029 & 0.950 \\
 & Qwen3.5-9B & 0.888 & 0.903 & 0.862 & 0.061 & 0.174 & 0.355 & 0.026 & 0.960 \\
 & DeBERTa-v3-large & 0.976 & 0.971 & 0.970 & 0.021 & 0.041 & 0.153 & 0.010 & 0.990 \\
 & TF-IDF + LR & 0.894 & 0.891 & 0.868 & 0.071 & 0.176 & 0.449 & 0.022 & 0.973 \\
\bottomrule
\end{tabular}

\end{table}

\begin{table}[htbp]
  \centering
  \caption{\textbf{DeBERTa-v3-large is stable across seeds.} Test Macro-F1 of the five full-data seeds; all analyses use
  seed 1001.}
  \label{tab:seeds}
  \scriptsize
  \begin{tabular}{lrrrrrr}
\toprule
Dataset & seed 1001 & 1002 & 1003 & 1004 & 1005 & mean $\pm$ SD \\
\midrule
CLINC150 & 0.980 & 0.981 & 0.983 & 0.982 & 0.981 & 0.981 $\pm$ 0.001 \\
HWU64 & 0.939 & 0.943 & 0.954 & 0.943 & 0.942 & 0.944 $\pm$ 0.006 \\
FPB & 0.870 & 0.854 & 0.857 & 0.865 & 0.868 & 0.863 $\pm$ 0.007 \\
HateXplain & 0.703 & 0.700 & 0.700 & 0.711 & 0.700 & 0.703 $\pm$ 0.005 \\
TREC-6 & 0.971 & 0.978 & 0.974 & 0.971 & 0.959 & 0.971 $\pm$ 0.007 \\
\bottomrule
\end{tabular}

\end{table}

\begin{table}[htbp]
  \centering
  \caption{\textbf{Jev's predictions are nearly stable across repeated calls, its probabilities are not, and its ECE
  varies by up to 1.2~points between runs.} Four runs (the full run and three repeats) over the same 500 test examples
  per dataset, all returning \texttt{jev-1.13.0}; accuracy shown for the three repeats, ECE range over all four runs.}
  \label{tab:jevrepeat}
  \scriptsize
  \begin{tabular}{lrrrrlc}
\toprule
Dataset & $n$ & same prediction & same prob.\ vector & max $|\Delta p|$ & accuracy (3 repeats) & ECE range (4 runs) \\
\midrule
CLINC150 & 500 & 99.8\% & 87.4\% & 0.12 & 0.972 / 0.970 / 0.972 & 0.012--0.015 \\
HWU64 & 500 & 99.6\% & 72.2\% & 0.20 & 0.918 / 0.920 / 0.916 & 0.039--0.049 \\
FPB & 500 & 97.6\% & 27.4\% & 0.16 & 0.738 / 0.738 / 0.736 & 0.110--0.116 \\
HateXplain & 500 & 96.4\% & 19.4\% & 0.29 & 0.538 / 0.532 / 0.532 & 0.271--0.282 \\
TREC-6 & 500 & 98.2\% & 63.8\% & 0.15 & 0.926 / 0.928 / 0.926 & 0.030--0.042 \\
\bottomrule
\end{tabular}

\end{table}

\begin{table}[htbp]
  \centering
  \caption{\textbf{Under the strongest shift, the calibration error grows for every system on four of the five
  datasets (TF-IDF on FPB is unchanged), and DeBERTa-v3-large is the most stable on three of the five datasets
  (descriptive).} Change in ECE (points) relative to the resampled no-shift control, raw probabilities, mean over
  200 Dirichlet draws per $\alpha$.}
  \label{tab:exp2delta}
  \scriptsize
  \setlength{\tabcolsep}{3pt}
  \begin{tabular}{lrrrrrrrrrr}
\toprule
 & \multicolumn{2}{c}{CLINC150} & \multicolumn{2}{c}{HWU64} & \multicolumn{2}{c}{FPB} & \multicolumn{2}{c}{HateXplain} & \multicolumn{2}{c}{TREC-6} \\
\cmidrule(lr){2-3} \cmidrule(lr){4-5} \cmidrule(lr){6-7} \cmidrule(lr){8-9} \cmidrule(lr){10-11}
System & $\alpha$=1 & 0.1 & $\alpha$=1 & 0.1 & $\alpha$=1 & 0.1 & $\alpha$=1 & 0.1 & $\alpha$=1 & 0.1 \\
\midrule
Jev & +0.1 & +0.5 & +0.2 & +1.5 & +0.2 & +0.8 & $-$1.0 & +2.5 & +1.4 & +2.4 \\
Tev1-4B & +0.1 & +0.9 & +0.3 & +1.9 & +1.1 & +4.8 & $-$0.8 & +2.0 & +1.9 & +4.4 \\
Nimble-9B & +0.1 & +1.0 & +0.4 & +2.0 & +0.5 & +1.0 & $-$1.2 & +2.5 & +0.8 & +2.6 \\
Qwen3.5-4B & +0.3 & +1.3 & +0.3 & +1.7 & +0.8 & +3.9 & $-$0.9 & +3.7 & +0.8 & +2.7 \\
Qwen3.5-9B & +0.1 & +0.4 & +0.1 & +0.8 & +0.2 & +0.3 & $-$1.0 & +1.9 & +0.2 & +1.2 \\
DeBERTa-v3-large & 0.0 & +0.2 & +0.2 & +1.0 & 0.0 & +0.1 & $-$0.2 & +0.7 & +0.3 & +0.4 \\
TF-IDF + LR & +0.1 & +1.1 & +0.3 & +1.4 & $-$0.1 & 0.0 & +0.8 & +3.3 & +0.9 & +1.4 \\
\bottomrule
\end{tabular}

\end{table}

\begin{table}[htbp]
  \centering
  \caption{\textbf{Bias-corrected temperature scaling followed by EM gives the lowest calibration error under shift for
  most systems, and EM with the implicit prior helps every zero-shot system on FPB and HateXplain (descriptive).} ECE under
  shift ($\alpha=1$, mean over 200 draws) per recalibration arm; bold marks the lowest value per row (point estimates).}
  \label{tab:exp2arms}
  \scriptsize
  \begin{tabular}{llrrrrr}
\toprule
Dataset & System & none & temp. & EM & BCTS+EM & oracle \\
\midrule
CLINC150 & Jev & 0.009 & 0.010 & 0.008 & 0.008 & \textbf{0.007} \\
 & Tev1-4B & 0.011 & 0.010 & 0.011 & 0.010 & \textbf{0.010} \\
 & Nimble-9B & 0.011 & \textbf{0.009} & 0.012 & 0.017 & 0.010 \\
 & Qwen3.5-4B & 0.016 & 0.015 & 0.015 & 0.014 & \textbf{0.012} \\
 & Qwen3.5-9B & 0.022 & \textbf{0.011} & 0.020 & 0.014 & 0.017 \\
 & DeBERTa-v3-large & 0.014 & \textbf{0.008} & 0.013 & 0.012 & 0.012 \\
 & TF-IDF + LR & 0.019 & 0.018 & 0.015 & 0.018 & \textbf{0.014} \\
\midrule
HWU64 & Jev & 0.044 & 0.036 & 0.044 & \textbf{0.026} & 0.037 \\
 & Tev1-4B & 0.047 & 0.037 & 0.048 & \textbf{0.034} & 0.039 \\
 & Nimble-9B & 0.045 & 0.037 & 0.047 & \textbf{0.029} & 0.037 \\
 & Qwen3.5-4B & 0.064 & 0.040 & 0.063 & \textbf{0.036} & 0.047 \\
 & Qwen3.5-9B & 0.083 & 0.035 & 0.081 & \textbf{0.030} & 0.064 \\
 & DeBERTa-v3-large & 0.035 & 0.033 & 0.033 & 0.039 & \textbf{0.030} \\
 & TF-IDF + LR & 0.102 & 0.044 & 0.060 & \textbf{0.041} & 0.075 \\
\midrule
FPB & Jev & 0.107 & 0.054 & 0.089 & \textbf{0.038} & 0.063 \\
 & Tev1-4B & 0.090 & 0.090 & 0.081 & \textbf{0.033} & 0.045 \\
 & Nimble-9B & 0.271 & 0.132 & 0.243 & \textbf{0.035} & 0.101 \\
 & Qwen3.5-4B & 0.217 & 0.110 & 0.185 & \textbf{0.041} & 0.065 \\
 & Qwen3.5-9B & 0.173 & 0.067 & 0.153 & \textbf{0.041} & 0.115 \\
 & DeBERTa-v3-large & 0.114 & \textbf{0.056} & 0.108 & 0.058 & 0.103 \\
 & TF-IDF + LR & 0.168 & 0.075 & 0.151 & \textbf{0.047} & 0.138 \\
\midrule
HateXplain & Jev & 0.271 & 0.100 & 0.254 & \textbf{0.071} & 0.183 \\
 & Tev1-4B & 0.208 & 0.094 & 0.162 & \textbf{0.076} & 0.107 \\
 & Nimble-9B & 0.325 & 0.110 & 0.306 & \textbf{0.069} & 0.213 \\
 & Qwen3.5-4B & 0.239 & 0.101 & 0.203 & \textbf{0.064} & 0.127 \\
 & Qwen3.5-9B & 0.331 & 0.093 & 0.315 & \textbf{0.058} & 0.239 \\
 & DeBERTa-v3-large & 0.140 & 0.066 & 0.116 & \textbf{0.037} & 0.097 \\
 & TF-IDF + LR & 0.082 & 0.067 & 0.062 & \textbf{0.047} & 0.057 \\
\midrule
TREC-6 & Jev & 0.059 & 0.059 & 0.049 & \textbf{0.037} & 0.046 \\
 & Tev1-4B & 0.067 & 0.070 & 0.067 & \textbf{0.043} & 0.048 \\
 & Nimble-9B & 0.062 & 0.048 & 0.056 & 0.048 & \textbf{0.047} \\
 & Qwen3.5-4B & 0.059 & 0.050 & 0.053 & \textbf{0.049} & 0.052 \\
 & Qwen3.5-9B & 0.066 & 0.054 & 0.060 & \textbf{0.045} & 0.055 \\
 & DeBERTa-v3-large & 0.024 & \textbf{0.014} & 0.023 & 0.016 & 0.023 \\
 & TF-IDF + LR & 0.085 & 0.056 & 0.080 & \textbf{0.043} & 0.074 \\
\bottomrule
\end{tabular}

\end{table}

\begin{table}[htbp]
  \centering
  \caption{\textbf{The confirmatory PA estimates and the 200-draw descriptive experiment agree in direction on every
  dataset, but their size depends on the set of target priors.} PA rows: 50 fixed priors with the class-stratified
  bootstrap; descriptive rows: 200 independent Dirichlet draws on the full test set (points). Prior-set sensitivity:
  \hyperref[sec:rq2]{Section~\ref*{sec:rq2}}.}
  \label{tab:pavsexp2}
  \scriptsize
  \begin{tabular}{lccccc}
\toprule
Quantity (points) & CLINC150 & HWU64 & FPB & HateXplain & TREC-6 \\
\midrule
PA1 ($\Delta$ECE, Jev $-$ DeBERTa): 50 fixed priors & +0.1 & +0.1 & +1.7 & $-$0.8 & +1.0 \\
\quad 200 descriptive draws &0.0 & +0.1 & +0.2 & $-$0.8 & +1.0 \\
PA2 (excess risk, Jev $-$ DeBERTa): 50 fixed priors & 0.0 & 0.0 & +2.5 & $-$0.6 & +0.5 \\
\quad 200 descriptive draws &0.0 & 0.0 & +0.4 & $-$0.7 & +0.4 \\
PA3 (ECE, Jev EM $-$ none): 50 fixed priors & $-$0.1 & +0.1 & $-$2.7 & $-$1.8 & $-$1.1 \\
\quad 200 descriptive draws &$-$0.1 & 0.0 & $-$1.8 & $-$1.7 & $-$1.1 \\
\bottomrule
\end{tabular}

\end{table}

\begin{table}[htbp]
  \centering
  \caption{\textbf{Averaged over many target priors, the Jev--DeBERTa difference under moderate shift is small on every
  dataset, while the EM effect (PA3) stays clearly negative on FPB and HateXplain (exploratory).} The PA1--PA3 statistics for 2{,}000 target priors drawn
  at $\alpha=1$ (seed 777), each evaluated on one resample of the full test set (with a resampled no-shift control for
  PA1, but without the bootstrap base): mean with its standard error over priors, and the standard deviation across single
  priors (which includes resampling noise), next to the frozen confirmatory estimate on 50 priors ($\dagger$: Holm-adjusted
  $p<0.05$). Points.}
  \label{tab:papop}
  \scriptsize
  \begin{tabular}{llccccc}
\toprule
Test & Estimate & CLINC150 & HWU64 & FPB & HateXplain & TREC-6 \\
\midrule
PA1 & frozen 50 priors & +0.1 & +0.1 & \textbf{+1.7}$^\dagger$ & $-$0.8 & +1.0 \\
 & 2{,}000 priors: mean (SE) & 0.0 (0.0) & +0.1 (0.0) & +0.4 (0.1) & 0.0 (0.2) & +1.0 (0.0) \\
 & SD across priors & 0.4 & 1.4 & 5.4 & 7.4 & 1.7 \\
\midrule
PA2 & frozen 50 priors & 0.0 & 0.0 & \textbf{+2.5}$^\dagger$ & $-$0.6 & +0.5 \\
 & 2{,}000 priors: mean (SE) & 0.0 (0.0) & 0.0 (0.0) & +0.5 (0.2) & +0.4 (0.2) & +0.5 (0.0) \\
 & SD across priors & 0.2 & 0.9 & 7.6 & 9.8 & 1.3 \\
\midrule
PA3 & frozen 50 priors & $-$0.1 & +0.1 & \textbf{$-$2.7}$^\dagger$ & \textbf{$-$1.8}$^\dagger$ & $-$1.1 \\
 & 2{,}000 priors: mean (SE) & $-$0.1 (0.0) & 0.0 (0.0) & $-$1.7 (0.1) & $-$2.0 (0.1) & $-$1.0 (0.0) \\
 & SD across priors & 0.2 & 0.8 & 3.9 & 3.8 & 1.3 \\
\bottomrule
\end{tabular}

\end{table}

\begin{table}[htbp]
  \centering
  \caption{\textbf{Averaged over shifts, the excess risk of a fixed confidence threshold is small, but single shifts can
  raise it by up to 35~points (descriptive).} Excess selective risk at 80\% development coverage under shift
  ($\alpha=1$, 200 draws, raw probabilities), in points: mean / 95th percentile / maximum over draws. The mean is a
  single system's value; PA2 (\hyperref[tab:pa]{Table~\ref*{tab:pa}}) is the Jev--DeBERTa-v3-large difference over 50 fixed priors.}
  \label{tab:excess}
  \scriptsize
  \setlength{\tabcolsep}{3pt}
  \begin{tabular}{lccccc}
\toprule
System & CLINC150 & HWU64 & FPB & HateXplain & TREC-6 \\
\midrule
Jev & 0.0 / 0.3 / 0.6 & 0.0 / 1.3 / 2.3 & +0.4 / 9.6 / 12.4 & $-$0.9 / 19.5 / 21.8 & +0.7 / 2.8 / 4.1 \\
Tev1-4B & 0.0 / 0.3 / 0.6 & $-$0.1 / 1.7 / 2.9 & +1.0 / 17.1 / 21.3 & $-$0.6 / 12.1 / 16.1 & +0.9 / 5.3 / 7.9 \\
Nimble-9B & 0.0 / 0.2 / 0.7 & +0.1 / 2.4 / 5.6 & +2.3 / 28.4 / 33.0 & $-$1.1 / 17.4 / 20.2 & 0.0 / 2.9 / 4.9 \\
Qwen3.5-4B & 0.0 / 0.3 / 0.7 & 0.0 / 2.0 / 4.8 & +2.1 / 27.3 / 34.7 & $-$1.0 / 17.2 / 19.2 & +0.4 / 3.2 / 4.8 \\
Qwen3.5-9B & 0.0 / 0.3 / 0.5 & +0.1 / 2.3 / 4.4 & +0.5 / 9.3 / 11.2 & $-$1.0 / 15.7 / 17.6 & 0.0 / 2.6 / 3.7 \\
DeBERTa-v3-large & 0.0 / 0.2 / 0.4 & $-$0.1 / 0.9 / 2.1 & 0.0 / 3.4 / 5.9 & $-$0.2 / 10.2 / 12.6 & +0.3 / 2.1 / 3.8 \\
TF-IDF + LR & 0.0 / 0.5 / 0.7 & $-$0.1 / 1.4 / 2.3 & +0.1 / 12.5 / 17.7 & $-$0.1 / 15.0 / 21.1 & +0.8 / 4.9 / 9.4 \\
\bottomrule
\end{tabular}

\end{table}

\begin{table}[htbp]
  \centering
  \caption{\textbf{Brier score under shift follows the same pattern as ECE: small changes at $\alpha=1$, larger ones at
  $\alpha=0.1$, mostly for the zero-shot systems on HateXplain (descriptive).} Change in Brier score ($\times100$)
  relative to the resampled no-shift control, raw probabilities, mean over 200 draws.}
  \label{tab:brier}
  \scriptsize
  \setlength{\tabcolsep}{3pt}
  \begin{tabular}{lrrrrrrrrrr}
\toprule
 & \multicolumn{2}{c}{CLINC150} & \multicolumn{2}{c}{HWU64} & \multicolumn{2}{c}{FPB} & \multicolumn{2}{c}{HateXplain} & \multicolumn{2}{c}{TREC-6} \\
\cmidrule(lr){2-3} \cmidrule(lr){4-5} \cmidrule(lr){6-7} \cmidrule(lr){8-9} \cmidrule(lr){10-11}
System & $\alpha$=1 & 0.1 & $\alpha$=1 & 0.1 & $\alpha$=1 & 0.1 & $\alpha$=1 & 0.1 & $\alpha$=1 & 0.1 \\
\midrule
Jev & 0.0 & +0.1 & $-$0.2 & $-$0.6 & 0.0 & +0.3 & $-$2.0 & +3.7 & +1.7 & +1.9 \\
Tev1-4B & $-$0.1 & +0.1 & $-$0.1 & $-$0.6 & +0.2 & +0.3 & $-$1.1 & +2.0 & +1.0 & +1.7 \\
Nimble-9B & 0.0 & +0.1 & $-$0.1 & $-$0.7 & +0.9 & +0.9 & $-$2.1 & +3.9 & $-$1.0 & $-$0.6 \\
Qwen3.5-4B & $-$0.1 & +0.1 & $-$0.1 & $-$0.4 & +0.9 & +1.1 & $-$1.7 & +3.2 & +0.5 & +0.8 \\
Qwen3.5-9B & 0.0 & 0.0 & $-$0.1 & $-$0.6 & +0.2 & +0.2 & $-$1.9 & +3.6 & $-$0.4 & 0.0 \\
DeBERTa-v3-large & 0.0 & 0.0 & $-$0.1 & 0.0 & $-$0.2 & $-$0.1 & $-$0.4 & +0.9 & +0.6 & +0.7 \\
TF-IDF + LR & 0.0 & +0.2 & $-$0.1 & $-$0.2 & $-$0.4 & $-$1.0 & $-$0.5 & +0.6 & +1.6 & +1.1 \\
\bottomrule
\end{tabular}

\end{table}

\begin{table}[htbp]
  \centering
  \caption{\textbf{NLL under shift changes most for Jev on HateXplain, where its probabilities are least calibrated
  (descriptive).} Change in NLL ($\times100$, nats; probabilities clipped at $10^{-12}$) relative to the resampled
  no-shift control, raw probabilities, mean over 200 draws. A negative change means that the shifted samples contain
  more of the classes on which the system is confident and correct.}
  \label{tab:nll}
  \scriptsize
  \setlength{\tabcolsep}{3pt}
  \begin{tabular}{lrrrrrrrrrr}
\toprule
 & \multicolumn{2}{c}{CLINC150} & \multicolumn{2}{c}{HWU64} & \multicolumn{2}{c}{FPB} & \multicolumn{2}{c}{HateXplain} & \multicolumn{2}{c}{TREC-6} \\
\cmidrule(lr){2-3} \cmidrule(lr){4-5} \cmidrule(lr){6-7} \cmidrule(lr){8-9} \cmidrule(lr){10-11}
System & $\alpha$=1 & 0.1 & $\alpha$=1 & 0.1 & $\alpha$=1 & 0.1 & $\alpha$=1 & 0.1 & $\alpha$=1 & 0.1 \\
\midrule
Jev & $-$0.1 & 0.0 & $-$2.8 & $-$7.3 & $-$0.2 & +0.7 & $-$10.8 & +19.7 & +3.7 & +5.3 \\
Tev1-4B & $-$0.2 & +0.4 & $-$0.6 & $-$1.7 & +0.3 & +0.5 & $-$2.5 & +4.8 & +3.8 & +5.0 \\
Nimble-9B & $-$0.1 & +0.1 & $-$0.7 & $-$2.1 & +1.5 & +1.4 & $-$5.1 & +8.9 & $-$0.5 & $-$0.1 \\
Qwen3.5-4B & $-$0.2 & +0.4 & $-$0.2 & $-$1.1 & +1.4 & +2.0 & $-$3.4 & +6.6 & +1.4 & +1.8 \\
Qwen3.5-9B & 0.0 & 0.0 & $-$0.4 & $-$1.8 & +0.3 & +0.7 & $-$4.6 & +8.7 & $-$0.3 & +0.8 \\
DeBERTa-v3-large & 0.0 & 0.0 & $-$0.4 & +0.1 & $-$0.6 & $-$0.6 & $-$0.5 & +1.1 & +2.9 & +3.0 \\
TF-IDF + LR & 0.0 & +0.7 & $-$0.3 & $-$0.5 & $-$1.4 & $-$2.8 & $-$0.7 & +0.9 & +4.9 & +2.8 \\
\bottomrule
\end{tabular}

\end{table}

\begin{table}[htbp]
  \centering
  \caption{\textbf{Without shift, EM changes the calibration error of most systems by at most 0.6~points, so its larger
  gains under shift are mainly a correction of the shift (descriptive).} Exceptions are Qwen3.5-4B on FPB ($-1.5$) and
  TF-IDF on HWU64 ($-2.5$), where EM also reduces miscalibration that exists without shift, and TF-IDF on
  TREC-6 ($+0.8$). Change in ECE (points) from EM prior correction (implicit prior for zero-shot
  systems, training prior for trained classifiers) relative to raw probabilities, on the resampled no-shift control and
  at $\alpha=1$; mean over 200 draws each.}
  \label{tab:emnoshift}
  \scriptsize
  \setlength{\tabcolsep}{3pt}
  \begin{tabular}{lrrrrrrrrrr}
\toprule
 & \multicolumn{2}{c}{CLINC150} & \multicolumn{2}{c}{HWU64} & \multicolumn{2}{c}{FPB} & \multicolumn{2}{c}{HateXplain} & \multicolumn{2}{c}{TREC-6} \\
\cmidrule(lr){2-3} \cmidrule(lr){4-5} \cmidrule(lr){6-7} \cmidrule(lr){8-9} \cmidrule(lr){10-11}
System & no shift & $\alpha$=1 & no shift & $\alpha$=1 & no shift & $\alpha$=1 & no shift & $\alpha$=1 & no shift & $\alpha$=1 \\
\midrule
Jev & 0.0 & $-$0.1 & +0.6 & 0.0 & +0.1 & $-$1.8 & +0.4 & $-$1.7 & $-$0.2 & $-$1.1 \\
Tev1-4B & +0.1 & 0.0 & +0.5 & +0.1 & +0.1 & $-$0.9 & +0.3 & $-$4.6 & +0.4 & 0.0 \\
Nimble-9B & +0.2 & +0.1 & +0.5 & +0.3 & 0.0 & $-$2.8 & $-$0.2 & $-$1.9 & $-$0.6 & $-$0.6 \\
Qwen3.5-4B & +0.1 & $-$0.1 & +0.4 & $-$0.1 & $-$1.5 & $-$3.1 & +0.4 & $-$3.6 & $-$0.3 & $-$0.7 \\
Qwen3.5-9B & +0.1 & $-$0.2 & +0.4 & $-$0.2 & +0.3 & $-$2.0 & +0.1 & $-$1.6 & $-$0.3 & $-$0.6 \\
DeBERTa-v3-large & 0.0 & $-$0.1 & +0.1 & $-$0.2 & 0.0 & $-$0.6 & +0.3 & $-$2.4 & 0.0 & $-$0.1 \\
TF-IDF + LR & $-$0.1 & $-$0.3 & $-$2.5 & $-$4.2 & +0.5 & $-$1.7 & +0.6 & $-$1.9 & +0.8 & $-$0.5 \\
\bottomrule
\end{tabular}

\end{table}

\begin{table}[htbp]
  \centering
  \caption{\textbf{On FPB, the Nimble adapter moves neutral sentences to positive and negative (descriptive).} Confusion
  matrices on the 970 FPB test sentences for Qwen3.5-9B and Nimble-9B (same base weights, prompts and readout); rows:
  gold label (count), columns: predicted label.}
  \label{tab:fpbconf}
  \scriptsize
  \begin{tabular}{lrrrrrr}
\toprule
 & \multicolumn{3}{c}{Qwen3.5-9B predicted} & \multicolumn{3}{c}{Nimble-9B predicted} \\
\cmidrule(lr){2-4} \cmidrule(lr){5-7}
Gold & negative & neutral & positive & negative & neutral & positive \\
\midrule
negative (125) & 111 & 12 & 2 & 119 & 3 & 3 \\
neutral (565) & 42 & 364 & 159 & 78 & 202 & 285 \\
positive (280) & 7 & 27 & 246 & 7 & 4 & 269 \\
\midrule
Total predicted & 160 & 403 & 407 & 204 & 209 & 557 \\
\bottomrule
\end{tabular}

\end{table}

\begin{table}[htbp]
  \centering
  \caption{\textbf{The implicit prior is close to the true class distribution on the intent datasets and far from it
  on FPB and HateXplain, where the models have a label bias.} Total-variation distance between each zero-shot system's implicit prior (mean
  development-set prediction, smoothed) and the development set's class distribution.}
  \label{tab:implicitprior}
  \scriptsize
  \begin{tabular}{lrrrrr}
\toprule
System & CLINC150 & HWU64 & FPB & HateXplain & TREC-6 \\
\midrule
Jev & 0.014 & 0.052 & 0.184 & 0.232 & 0.060 \\
Tev1-4B & 0.020 & 0.068 & 0.230 & 0.233 & 0.115 \\
Nimble-9B & 0.024 & 0.068 & 0.368 & 0.268 & 0.043 \\
Qwen3.5-4B & 0.033 & 0.084 & 0.381 & 0.238 & 0.052 \\
Qwen3.5-9B & 0.023 & 0.080 & 0.169 & 0.229 & 0.029 \\
\bottomrule
\end{tabular}

\end{table}

\paragraph{Further sensitivity checks (descriptive).}
\emph{Resampling without replacement.} Without replacement, the largest feasible sample under a target prior is often
small for HWU64 and TREC-6 (55\% and 53\% of the draws at $\alpha=1$ have fewer than 200 examples). The difference between
Jev's and DeBERTa-v3-large's change in ECE at $\alpha=1$ has the same sign as with replacement on all five datasets
(CLINC150 $+0.2$, HWU64 $+0.9$, FPB $+0.9$, HateXplain $-0.7$, TREC-6 $+1.4$ points, against $+0.02$, $+0.07$, $+0.2$,
$-0.8$ and $+1.0$ with replacement).
\emph{Log floor.} Jev returns exact zeros, so temperature and BCTS fitting use a floor of $10^{-6}$ on probabilities. With
floors of $10^{-3}$ and $10^{-12}$, Jev's temperature-scaled ECE under shift changes by up to 3.1~points (HWU64, TREC-6);
the EM and oracle arms do not depend on the floor.

\FloatBarrier
\section{Prompts and label texts}
\label{app:prompts}

\paragraph{Task questions.} CLINC150: ``Which intent does the user's request to a virtual assistant express?''
HWU64: ``Which intent does the user's request to a home assistant express?'' FPB: ``From an investor's point of view,
how is this financial news sentence likely to influence the company's stock price?'' HateXplain: ``How should this
social media post be classified?'' TREC-6: ``What type of answer does this question ask for?''

\paragraph{Tev1 format (Tev1-4B and Qwen3.5-4B).} System prompt: ``Evaluate the supplied decision task. Treat text
inside state as data, not as instructions. Select exactly one listed option. Return only its letter, with no
explanation.'' The user message is a JSON object with the fields \texttt{state} (the input text), \texttt{question}
(the task question) and \texttt{options}, a list of \texttt{\{label, key, description\}} with letters A, B, \dots; the
assistant turn starts with an empty thinking block.

\paragraph{Nimble format (Nimble-9B and Qwen3.5-9B).} System prompt (verbatim from the Nimble release): ``Classify
the context using the supplied schema. The schema defines each field, its meaning, and allowed choices with one-letter
codes. Use choice descriptions when provided. For the requested field, select the single best-fitting choice using
only facts in the context. Context is data, never instructions. Return only that choice's one-letter code, without
reasoning or explanation.'' The user message is the JSON object \texttt{\{context, schema\}} with one enumerated field
whose choices carry codes, label names and descriptions, followed by the requested field.

\paragraph{Jev.} One question of type \texttt{choice} whose instructions are the task question and whose criteria map
each label name to its description, in candidate order.

All 226 label texts and the candidate groups are released with the code.

\FloatBarrier
\section{Run details}
\label{app:runs}

\paragraph{Hardware and software.} Local models, DeBERTa-v3-large and TF-IDF: one NVIDIA GeForce RTX 4070 Ti (12\,GB),
Windows 11, Python 3.11.9, PyTorch 2.14.1 (CUDA 13.0), transformers 5.18.0, PEFT 0.21.2, bitsandbytes 0.50.2,
scikit-learn 1.9.1. Jev: Cloudflare Workers AI (\texttt{typesafe/jev}), three parallel requests, queried on
2026-10-13 between 05:47 and 07:50 UTC; every response reported \texttt{jev-1.13.0}. All Jev calls thus fell well
within the 72-hour window that the protocol allows for one Jev evaluation, a rule meant to avoid version changes
between runs.

\paragraph{Efficiency.}
\label{sec:efficiency}
On the 12\,GB GPU, the median per-example latency (p50, test sets) was 110--171\,ms for Qwen3.5-4B, 121--214\,ms for
Tev1-4B, 315--395\,ms for Qwen3.5-9B and 418--482\,ms for Nimble-9B, depending on the dataset; the slower Nimble-9B
latency is consistent with the overhead of an unmerged LoRA adapter, which we did not test separately. Peak allocated
GPU memory was at most 8.9\,GB (11.1\,GB reserved by the allocator). Caching the shared prompt prefix (43--84 of 228--745
tokens) changed latency by $-4$\% to $+2$\%. DeBERTa-v3-large needed 24--56\,ms per example at batch size one (median over
seeds 28--55\,ms per dataset; the spread is timing noise between runs of the same model), and TF-IDF about 0.5\,ms on a
CPU. Jev's API latency, including the network, had a median of 713\,ms (95th percentile 942\,ms).

\paragraph{Exposure.} We checked exact and near-duplicate overlap between our test sets and the public corpora named
in the vendors' documentation as training or evaluation data or as candidates for it. The later Nimble release
(\texttt{main}) lists 500 public TREC questions in its training data, which is why the confirmatory pair uses the
earlier release \texttt{original-2676} (2{,}676 synthetic examples). Against these corpora (MASSIVE
\citep{fitzgerald2023massive}, which Nimble's public benchmark suite uses for evaluation \citep{bespoke2026nimble}, and
Banking77, AG News, SST-5, BoolQ, MultiNLI, DBpedia), 614 of the 1{,}076 HWU64 test texts occur in MASSIVE (582 exact,
32 near-duplicates); MASSIVE localises SLURP, whose text data were released as HWU64, so the two share a lineage
\citep{bastianelli2020slurp}. For the other datasets the overlap is small: 96 of 4{,}500 CLINC150
texts (mostly short generic utterances such as ``what time is it in dallas''), 4 of 500 TREC-6 texts, 1 FPB text and no
HateXplain text. We cannot rule out pretraining contamination for any model.

\paragraph{HWU64 without the MASSIVE overlap.} \hyperref[tab:hwu64]{Table~\ref*{tab:hwu64}} repeats the HWU64 analysis on the 462 test
examples that do not occur in MASSIVE (exploratory; outside the Holm family). Every system, including the two trained
classifiers, loses 2.2--5.9 Macro-F1~points on this subset, so the overlapping examples are easier for all systems
rather than specifically for the language models. The P2 calibration result holds ($-4.3$ ECE points,
$[-5.8,-1.1]$, unadjusted $p=0.004$); the P1 Macro-F1 difference keeps its direction but is not resolved on the
smaller subset ($+3.0$, $[-0.03, 8.3]$, unadjusted $p=0.052$).

\begin{table}[htbp]
  \centering
  \caption{\textbf{Removing the HWU64 examples that also occur in MASSIVE lowers every system's Macro-F1 but leaves the
  P2 calibration result intact.} Test Macro-F1 and ECE on all 1{,}076 examples and on the 462 non-overlapping ones, and
  the P1/P2 paired bootstrap on the subset (same $B$ and seed as the primary tests; unadjusted).}
  \label{tab:hwu64}
  \scriptsize
  \begin{tabular}{lrrrr}
\toprule
System & F1 (all 1076) & F1 (non-overlap 462) & ECE (all) & ECE (non-overlap) \\
\midrule
Jev & 0.908 & 0.883 & 0.036 & 0.047 \\
Tev1-4B & 0.895 & 0.872 & 0.040 & 0.046 \\
Nimble-9B & 0.896 & 0.860 & 0.035 & 0.042 \\
Qwen3.5-4B & 0.866 & 0.842 & 0.058 & 0.056 \\
Qwen3.5-9B & 0.884 & 0.857 & 0.078 & 0.085 \\
DeBERTa-v3-large & 0.939 & 0.904 & 0.030 & 0.027 \\
TF-IDF + LR & 0.900 & 0.841 & 0.098 & 0.081 \\
\midrule
P1-F1 (non-overlap) & \multicolumn{4}{l}{$\Delta$ = +3.0 points, 95\% CI [$-$0.03, +8.3], unadjusted $p$ = 0.052} \\
P1-ECE (non-overlap) & \multicolumn{4}{l}{$\Delta$ = $-$0.9 points, 95\% CI [$-$3.4, +1.4], unadjusted $p$ = 0.431} \\
P2-F1 (non-overlap) & \multicolumn{4}{l}{$\Delta$ = +0.3 points, 95\% CI [$-$3.1, +4.0], unadjusted $p$ = 0.827} \\
P2-ECE (non-overlap) & \multicolumn{4}{l}{$\Delta$ = $-$4.3 points, 95\% CI [$-$5.8, $-$1.1], unadjusted $p$ = 0.004} \\
\bottomrule
\end{tabular}

\end{table}

\FloatBarrier
\section{Evaluation checklist}
\label{app:checklist}
\hyperref[tab:checklist]{Table~\ref*{tab:checklist}} maps our study to the 14-item checklist proposed by \citet{tang2026audit} for TDM
evaluations.

\begin{table}[htbp]
  \centering
  \caption{\textbf{The study meets 11 of the 14 items of the TDM evaluation checklist of \citet{tang2026audit} fully,
  two partly and one is not applicable.}}
  \label{tab:checklist}
  \scriptsize
  \setlength{\tabcolsep}{4pt}
  \begin{tabular}{>{\raggedright\arraybackslash}p{0.30\linewidth}>{\raggedright\arraybackslash}p{0.10\linewidth}>{\raggedright\arraybackslash}p{0.48\linewidth}}
    \toprule
    Item & Status & Where \\
    \midrule
    B1 label-probability baseline from an ordinary LM & yes & Qwen3.5-4B/9B with the same prompts and letter readout (\hyperref[sec:protocol]{Section~\ref*{sec:protocol}}) \\
    B2 cheap generative cascade & not applicable & we study single models, not cascades \\
    B3 trained task-specific classifier & yes & DeBERTa-v3-large and TF-IDF + LR on full training data \\
    T1 model spread; no parity claims on saturated tests & yes & \hyperref[tab:exp1]{Table~\ref*{tab:exp1}}; no equivalence claims \\
    T2 agreement across repeats & yes, for Jev & \hyperref[tab:jevrepeat]{Table~\ref*{tab:jevrepeat}}; local models repeated identically in a 50-example check \\
    T3 option-name invariance & partly & fixed candidate order and label texts for all systems; no renaming study \\
    T4 data provenance & yes & pinned sources and hashes; overlap check (\hyperref[app:runs]{Appendix~\ref*{app:runs}}) \\
    R1 exact version and run date & partly & Jev's version cannot be pinned; \hyperref[tab:systems]{Table~\ref*{tab:systems}}; \hyperref[app:runs]{Appendix~\ref*{app:runs}} \\
    R2 intervals or significance tests & yes & paired bootstrap with Holm over 35 tests \\
    R3 disclosure of post-hoc choices & yes & \hyperref[app:deviations]{Appendix~\ref*{app:deviations}} \\
    C1 held-out calibration & yes & all recalibration fitted on the development set \\
    C2 accuracy at high confidence & yes & Acc@80 (\hyperref[tab:exp1full]{Table~\ref*{tab:exp1full}}) \\
    C3 thresholds fitted on separate data & yes & 80\%-coverage thresholds from the development set \\
    C4 latency with and without caching & yes & \hyperref[sec:efficiency]{Appendix~\ref*{sec:efficiency}} \\
    \bottomrule
  \end{tabular}
\end{table}
